\documentclass[12pt]{article}

\usepackage{sbc-template}
\usepackage{graphicx,url}
\usepackage[utf8]{inputenc}
\usepackage[brazil]{babel}
\usepackage{booktabs}
\usepackage{multirow}
\usepackage{amsmath}

\sloppy

\title{Solução IDS Hierárquica para Detecção de Ciberataques em Ambientes IoT com MEC em Redes 6G}

\author{Luiz Henrique B. A. da Silva\inst{1}}

\address{Curso de Sistemas de Informação\\
  Universidade Federal de Pernambuco -- Recife -- PE -- Brasil\\
  \texttt{lhbas@cin.ufpe.br}
}

\begin{document}

\maketitle

%─────────────────────────────────────────────────────────────────────────────
\begin{abstract}
The exponential growth of Internet of Things (IoT) devices and the emergence
of sixth-generation (6G) networks introduce new cybersecurity challenges in
Multi-access Edge Computing (MEC) environments. This work proposes and
implements a two-phase hierarchical Intrusion Detection System (IDS) for
IoT/MEC environments. Phase 1 runs a binary LightGBM classifier on the edge
node, optimized by the F2-score; Phase 2 classifies the specific attack type
using a multiclass LightGBM. The data
pipeline uses PCAP captures from five CIC datasets converted by
\textit{netflower} — a tool developed for this project — into a 190~GB SQLite
database. In offline evaluation, Phase 1 achieved 80\% accuracy and 96\%
attack recall (F2 = 0.886). In live validation on VIM~4, the binary IDS
achieved 85.3\% accuracy, 99.6\% precision, 78.3\% recall and 87.7\%
F1-score, with FPR of 0.7\%. The multiclass IDS added attack-type
classification at 81.3\% accuracy with F1 of 0.884 for DoS and 0.756 for
reconnaissance, at the cost of 52~MB additional RAM. Results demonstrate the
viability of the hierarchical approach within the resource budget of an ARM
edge node.
\end{abstract}

%─────────────────────────────────────────────────────────────────────────────
\begin{resumo}
O crescimento exponencial de dispositivos da Internet das Coisas (IoT) e a
emergência das redes de sexta geração (6G) introduzem novos desafios de
segurança cibernética em ambientes de Computação na Borda Móvel (MEC).
Este trabalho propõe e implementa um Sistema de Detecção de Intrusão (IDS)
hierárquico de duas fases para ambientes IoT/MEC. A Fase~1 executa um
classificador binário LightGBM no nó de borda, otimizado pelo escore F2; a
Fase~2 classifica o tipo específico de ataque com LightGBM multiclasse. O pipeline utiliza PCAPs de
cinco datasets do CIC, convertidos pela ferramenta \textit{netflower} em banco
SQLite de $\sim$190~GB. Em avaliação offline local, a Fase~1 atingiu acurácia de
80\% e revocação de 96\% na classe de ataque (F2 = 0,886). Na validação ao
vivo no VIM~4, o IDS binário alcançou acurácia de 85,3\%, precisão de 99,6\%,
revocação de 78,3\% e F1 de 87,7\%, com FPR de 0,7\%. O IDS multiclasse
adicionou classificação de tipo com acurácia de 81,3\% (F1 de 0,884 para
DoS/DDoS e 0,756 para reconhecimento), ao custo de apenas 52~MB extras de
RAM. Os resultados confirmam a viabilidade da abordagem hierárquica dentro do
orçamento de recursos de um nó de borda ARM.
\end{resumo}

%─────────────────────────────────────────────────────────────────────────────
\section{Introdução}

A convergência entre redes 6G, IoT e MEC promete transformar a infraestrutura
de comunicação digital. As redes 6G preveem densidades de dispositivos IoT na
ordem de $10^6$ por km², com latências inferiores a 1~ms, tornando o MEC
componente crítico para viabilizar esses requisitos~\cite{you:21}. Contudo,
essa massificação expande igualmente a superfície de ataque: dispositivos IoT
frequentemente possuem poder computacional e energia limitados, inviabilizando
soluções de segurança tradicionais diretamente nos terminais.

Sistemas de Detecção de Intrusão (IDS) baseados em aprendizado de máquina,
executados nos nós de borda (VIMs), surgem como alternativa viável: modelos
treinados em dados históricos de tráfego classificam fluxos de rede em tempo
real com alta eficiência computacional. Entretanto, os IDS existentes
enfrentam limitações para ambientes IoT/MEC: \textit{(i)} conjuntos de dados
capturados em ambientes controlados geram desvio de distribuição no tráfego
real; \textit{(ii)} a variedade de tipos de ataque exige diferentes objetivos
de otimização em cada camada hierárquica.

Este trabalho apresenta um IDS hierárquico de duas fases que aborda essas
limitações, descrevendo o pipeline completo desde a coleta de PCAPs até a
implantação e validação em produção no VIM~4. As contribuições incluem:
\textit{(i)} pipeline de dados reprodutível baseado em SQLite para amostragem
balanceada em escala; \textit{(ii)} ferramenta \textit{netflower} (publicada
no PyPI) para conversão paralela de PCAPs em fluxos; \textit{(iii)} IDS
hierárquico com duas fases complementares; \textit{(iv)} validação ao vivo com
tráfego de ataque real e comparação quantitativa entre o pipeline binário e o
multiclasse, incluindo métricas de acurácia, precisão, revocação, F1,
FPR/FNR, uso de CPU, RAM e throughput de alertas.

O artigo está organizado da seguinte forma: a Seção~2 apresenta os trabalhos
relacionados e a fundamentação teórica; a Seção~3 descreve a metodologia
proposta; a Seção~4 apresenta os resultados da validação ao vivo; a Seção~5
conclui o trabalho; a Seção~6 lista os trabalhos futuros.

%─────────────────────────────────────────────────────────────────────────────
\section{Trabalhos Relacionados}

IDS podem ser classificados em baseados em assinatura e baseados em anomalia.
Os primeiros identificam padrões conhecidos com alta precisão, mas são
incapazes de detectar ameaças novas. Os baseados em anomalia aprendem um perfil
do comportamento normal e sinalizam desvios, sendo mais adequados para
ambientes dinâmicos~\cite{lazarevic:05}. A combinação das duas abordagens em
arquitetura hierárquica representa o estado da arte para ambientes com
restrições de recursos como nós de borda MEC~\cite{zolanvari:19}.

O \textbf{LightGBM}~\cite{ke:17} é um algoritmo de gradient boosting baseado
em árvores de decisão com crescimento folha a folha (\textit{leaf-wise}), que
proporciona convergência rápida e desempenho superior em conjuntos de alta
dimensionalidade. O \textbf{Optuna}~\cite{akiba:19} é um framework de
otimização automática de hiperparâmetros baseado em TPE (\textit{Tree-structured
Parzen Estimator}), com suporte a poda de tentativas pouco promissoras.

%─────────────────────────────────────────────────────────────────────────────
\section{Metodologia}

A motivação para a arquitetura hierárquica de duas fases parte das restrições
inerentes a ambientes IoT/MEC: \textit{(i)} a maioria do tráfego em qualquer
rede é benigna — aplicar um classificador pesado sobre todos os fluxos
desperdiçaria energia e capacidade no nó de borda; \textit{(ii)} diferentes etapas do
pipeline têm objetivos de otimização distintos — detecção binária privilegia
revocação (não deixar ataques passarem), enquanto classificação de tipo privilegia
precisão por classe. O pipeline progressivo resolve essas restrições: a Fase~1
filtra tráfego benigno com custo mínimo; apenas fluxos suspeitos avançam à
Fase~2 para classificação de tipo.

A arquitetura proposta distribui a carga computacional hierarquicamente:

\begin{enumerate}
  \item \textbf{Fase~1 — Classificador binário (nó de borda):} decide
    \textit{benigno vs.\ ataque} com mínimo custo computacional.
  \item \textbf{Fase~2 — Classificador multiclasse (borda/MEC):} identifica o
    tipo específico de ataque para os fluxos flagados pela Fase~1.
\end{enumerate}

\subsection{Pipeline de dados}

\textbf{Datasets.} Apenas datasets com PCAPs separados por categoria foram
utilizados — os rótulos derivam do nome das pastas (\texttt{benign/},
\texttt{ddos/}, \texttt{malware/} etc.), evitando inconsistências de esquema
entre fontes~\cite{sharafaldin:18}. Foram empregados cinco datasets do CIC
(Canadian Institute for Cybersecurity), descritos na Tabela~\ref{tab:datasets}.

\begin{table}[ht]
\centering
\caption{Datasets CIC utilizados no pipeline de dados}
\label{tab:datasets}
\begin{tabular}{@{}ll@{}}
\toprule
\textbf{Dataset} & \textbf{Foco principal} \\ \midrule
CIC APT IIoT 2024    & Ameaças persistentes avançadas (APT) em IIoT \\
CIC IoT Dataset 2023 & Tráfego multi-ataque em dispositivos IoT \\
CIC IoT DIAD 2024    & Detecção distribuída de ataques IoT \\
CIC IIoT 2025        & Intrusões em IoT industrial \\
BCCC NRC IoMT 2024   & Segurança em IoT médico (dispositivos de saúde) \\
\bottomrule
\end{tabular}
\end{table}

Além do tráfego benigno proveniente dos datasets, tráfego real capturado de uso
normal de computadores (navegação web, sessões SSH, atualizações de sistema) foi
integrado à classe benigna. O tráfego dos datasets CIC é gerado de forma
\textit{scripted} com padrões temporais regulares; o tráfego interativo real
apresenta características distintas — \textit{bursts} de SSH, \textit{keepalives}
HTTP, fluxos de curta duração sobre rede local rápida — que modelos treinados
exclusivamente nos datasets tendem a classificar erroneamente como ataque. A
mistura reduz o desvio de distribuição entre treino e implantação.

\textbf{Pipeline de construção do banco.} A construção do banco segue três
etapas reprodutíveis:

\begin{enumerate}
  \item \textbf{PCAP $\rightarrow$ CSV (netflower):} para cada pasta de categoria
    nos datasets (e.g.\ \texttt{benign/}, \texttt{ddos/}, \texttt{malware/}),
    todos os arquivos \texttt{.pcap}/\texttt{.pcapng} são convertidos em fluxos
    via \texttt{convert\_pcap\_to\_csv(n\_jobs=8)} — oito processos paralelos
    particionando os pacotes por chave bidirecional determinística. O rótulo de
    cada fluxo é derivado automaticamente do nome da pasta, eliminando rotulagem
    manual e inconsistências de esquema entre fontes. Os fluxos de todos os
    PCAPs de uma mesma pasta são concatenados em um único arquivo
    \texttt{merged\_<label>.csv}.
  \item \textbf{CSV $\rightarrow$ SQLite:} cada \texttt{merged\_*.csv} é
    inserido no banco em \textit{chunks} de 50.000 linhas, sem carregar o arquivo
    inteiro em memória. Ao final, cria-se um índice na coluna \texttt{label}
    para viabilizar consultas eficientes por classe. O resultado é um banco de
    $\sim$190~GB com compressão de $\sim$5$\times$ frente às PCAPs originais
    ($\sim$1~TB).
  \item \textbf{Amostragem balanceada:} durante o treinamento, utiliza-se
    \texttt{ORDER BY RANDOM() LIMIT n} por classe diretamente no SQLite,
    garantindo amostras uniformes sem varredura completa do banco nem carga de
    todo o dataset em memória.
\end{enumerate}

\textbf{netflower: projeto e razões da escolha.} O \textit{netflower} é uma
ferramenta desenvolvida para este projeto e publicada no PyPI. Ela extrai fluxos
de rede bidirecionais a partir de arquivos \texttt{.pcap}/\texttt{.pcapng} ou
de interfaces ao vivo, produzindo 82~features compatíveis com o conjunto do
CICFlowMeter. A ferramenta foi projetada especificamente para execução em
dispositivos de borda, reunindo quatro características relevantes: \textit{(i)}
\textbf{desempenho em ARM} — o parser \textit{dpkt} é 10--25$\times$ mais rápido
que Scapy em processadores ARM, e o algoritmo de Welford acumula estatísticas em
memória $O(1)$ por fluxo, sem armazenar listas de pacotes; \textit{(ii)}
\textbf{captura ao vivo sem extensão compilada} — libpcap é acessado via
\texttt{ctypes}, eliminando dependências de compilação no nó de borda; \textit{(iii)}
\textbf{conversão paralela de PCAPs} — os pacotes são particionados por chave
bidirecional determinística (garantindo que pacotes de um mesmo fluxo sempre vão
ao mesmo worker) entre \texttt{n\_jobs} processos paralelos, com fusão posterior
em um único CSV; \textit{(iv)} \textbf{emissão orientada a fluxo} — cada fluxo é
emitido apenas ao completar (TCP FIN/RST ou \textit{timeout}), garantindo features
completas em toda linha.

\subsection{Mapeamento de rótulos}

Após análise das distribuições das classes, aplicou-se
$\texttt{LABEL\_MAP} = \{\text{"ddos"} \rightarrow \text{"dos"}\}$ antes do
treinamento. A justificativa: ataques DDoS e DoS visam exaurir recursos; seus
fluxos de rede (alta taxa de pacotes, curta duração, anomalias de flags)
são estatisticamente similares o suficiente para que manter duas classes
separadas adicione ruído sem ganho discriminativo significativo. A unificação
reduz o espaço de classes de 9 para 8 (7 de ataque + 1 benigno) e simplifica
o problema da Fase~2.

Do ponto de vista do nó VIM~4 (alvo), a distinção entre DoS e DDoS é de natureza
\textit{externa} aos fluxos individuais: um ataque DoS provém de uma única fonte,
enquanto o DDoS usa múltiplas fontes distribuídas — mas os fluxos individuais
recebidos pelo alvo são estatisticamente indistinguíveis. A feature que capturaria
essa diferença (diversidade de IPs de origem) exige agregação \textit{entre} fluxos
e não está presente nas 82~features por-fluxo do \textit{netflower}. Manter as
duas classes separadas no treinamento produziria um classificador que aprenderia a
distingui-las por artefatos dos datasets, não por características do tráfego real
observado pelo nó de borda.

A Tabela~\ref{tab:taxonomy} apresenta a taxonomia resultante.

\begin{table}[ht]
\centering
\caption{Taxonomia de ataques utilizada no treinamento do sistema IDS}
\label{tab:taxonomy}
\begin{tabular}{@{}ll@{}}
\toprule
\textbf{Categoria} & \textbf{Descrição} \\ \midrule
\texttt{benign}     & Tráfego legítimo de rede \\
\texttt{dos}        & DoS e DDoS (unificados) — exaustão de recursos \\
\texttt{malware}    & Tráfego de malware e comunicação C2 \\
\texttt{bruteforce} & Ataques de força bruta (SSH, HTTP, Telnet) \\
\texttt{mitm}       & Ataques Man-in-the-Middle (envenenamento ARP) \\
\texttt{web}        & Ataques a aplicações web (scanning, fuzzing) \\
\texttt{spoofing}   & Falsificação de identidade de origem em pacotes IP \\
\texttt{recon}      & Varreduras e reconhecimento de rede (port scan) \\ \bottomrule
\end{tabular}
\end{table}

\subsection{Engenharia de características}

Partindo de 83 colunas brutas produzidas pelo \textit{netflower}, o pipeline
de limpeza aplicou os seguintes critérios de remoção:

\begin{itemize}
  \item \textbf{Identidade de fluxo} (\texttt{src\_ip}, \texttt{dst\_ip},
    \texttt{src\_port}, \texttt{dst\_port}, \texttt{timestamp}): colunas que
    identificam hosts/serviços específicos fazem o modelo memorizar endereços
    vistos no treino em vez de aprender padrões estatísticos de tráfego —
    qualquer host novo no ambiente de implantação produziria viés de previsão.

  \item \textbf{Sentinel de janela TCP} (\texttt{init\_fwd\_win\_byts},
    \texttt{init\_bwd\_win\_byts}): o \textit{netflower} escreve $-1$ quando o
    tamanho da janela do handshake TCP não está disponível. Esse valor $-1$ não
    ocorre no tráfego real; ao manter a feature, o modelo aprende a disparar
    sobre o sentinel, gerando falsos positivos em fluxos legítimos onde o campo
    é genuinamente ausente.

  \item \textbf{Bulk rates} (\texttt{fwd\_byts\_b\_avg}, etc.): zerados em
    virtualmente todo tráfego real de implantação (exigem heurísticas de
    detecção de bulk não aplicadas pelo \textit{netflower} em modo ao vivo),
    mas não-nulos nos datasets de treino — criando desvio treino/implantação.

  \item \textbf{Correlação alta} ($> 0{,}95$): 11 features redundantes.
  \item \textbf{Variância próxima de zero}: 3 features sem sinal.
\end{itemize}

O resultado foram \textbf{55 features numéricas} finais. As duas de maior
relevância pelo ganho do LightGBM foram \texttt{flow\_byts\_s} (taxa de bytes
por segundo) e \texttt{flow\_iat\_min} (tempo mínimo de chegada entre pacotes),
que capturam, respectivamente, a intensidade e o padrão temporal do tráfego.

\subsection{Treinamento das Duas Fases}

\subsubsection{Fase 1 — Classificador Binário}

\textbf{Configuração:} LightGBM binário, 615.317 amostras benignas + 615.317
de ataque (distribuídas uniformemente entre as oito classes do SQLite), divisão
80/20 estratificada. O limiar de decisão foi otimizado conjuntamente com os
hiperparâmetros via Optuna (20 tentativas, 1.800~s de \textit{timeout}),
usando F2-score como objetivo — que pondera revocação $4\times$ mais que
precisão, refletindo que, em um IDS (não IPS), um falso negativo é mais
custoso que um falso positivo.

\textbf{Resultados (conjunto de teste — 140.712 fluxos, limiar = 0,157):}

\begin{table}[ht]
\centering
\caption{Avaliação offline da Fase~1 (modelo 20260601, limiar ótimo = 0,157)}
\label{tab:phase1_offline}
\begin{tabular}{@{}lrrrr@{}}
\toprule
\textbf{Classe} & \textbf{Precisão} & \textbf{Revocação} & \textbf{F1} & \textbf{Suporte} \\ \midrule
attack  & 0,68 & \textbf{0,96} & 0,80 & 57.398 \\
benign  & \textbf{0,96} & 0,69 & 0,80 & 83.314 \\ \midrule
\textbf{Acurácia} & & & \textbf{0,80} & 140.712 \\
Macro avg & 0,82 & 0,83 & 0,80 & \\
\bottomrule
\end{tabular}
\end{table}

\textbf{F2-score} (objetivo Optuna): $\mathbf{0{,}886}$. O modelo prioriza
revocação de ataques aceitando menor precisão — na prática, aproximadamente
1 em 3 alertas seria falso positivo ao limiar ótimo. O limiar de
implantação (\textit{deployed}) foi definido em $0{,}9$ para minimizar falsos
positivos em produção, ao custo de menor revocação.

\subsubsection{Fase 2 — Classificador Multiclasse}

\textbf{Configuração:} LightGBM multiclasse, amostragem independente de até
615.317 linhas por classe diretamente do SQLite, \texttt{class\_weight=balanced}
para compensar a escassez de \textit{bruteforce} ($\sim$16~mil amostras).
Objetivo Optuna: macro F1, garantindo equidade entre classes.

\textbf{Resultados (conjunto de teste — 317.381 fluxos):}

\begin{table}[ht]
\centering
\caption{Avaliação offline da Fase~2 por classe (modelo 20260601)}
\label{tab:phase2_offline}
\begin{tabular}{@{}lrrrr@{}}
\toprule
\textbf{Classe} & \textbf{Precisão} & \textbf{Revocação} & \textbf{F1} & \textbf{Suporte} \\ \midrule
benign      & 0,83 & 0,69 & 0,76 & 83.313 \\
bruteforce  & 0,20 & 0,77 & 0,31 &  2.191 \\
dos         & \textbf{0,99} & \textbf{0,96} & \textbf{0,98} & 64.582 \\
malware     & 0,90 & 0,78 & 0,84 & 81.338 \\
mitm        & 0,42 & 0,72 & 0,53 & 13.823 \\
recon       & 0,83 & 0,68 & 0,75 & 47.915 \\
spoofing    & 0,28 & 0,73 & 0,41 & 12.289 \\
web         & 0,77 & 0,84 & 0,80 & 11.930 \\ \midrule
\textbf{Acurácia} & & & \textbf{0,78} & 317.381 \\
Macro F1    & & & \textbf{0,67} & \\
\bottomrule
\end{tabular}
\end{table}

%─────────────────────────────────────────────────────────────────────────────
\section{Resultados}

A validação foi realizada no VIM~4 (Khadas VIM4, ARM Cortex-A73 + A53
octa-core, 8~GB RAM), com tráfego de ataque gerado pelo
\texttt{attack\_orchestrator.py} a partir de um PC atacante (BRT, UTC$-$3)
e os dois pipelines IDS em execução simultânea (binário e multiclasse).

\subsection{Configuração experimental}

O experimento envolveu dois equipamentos conectados na mesma rede local
(Tabela~\ref{tab:hardware}). O PC pessoal foi responsável tanto pelo treinamento
dos modelos quanto pela geração dos ataques; o VIM~4 atuou como nó de borda MEC
executando os scripts IDS via SSH e monitorando a interface \texttt{eth0}.

\begin{table}[ht]
\centering
\caption{Especificações do hardware utilizado}
\label{tab:hardware}
\begin{tabular}{@{}lll@{}}
\toprule
\textbf{Componente} & \textbf{PC (treino / ataques)} & \textbf{VIM 4 (IDS)} \\ \midrule
SoC / CPU & AMD Ryzen 5 7600 (6c/12t) & Amlogic A311D2 \\
          &                            & 4$\times$ Cortex-A73 @ 2,2~GHz \\
          &                            & + 4$\times$ Cortex-A53 @ 1,8~GHz \\
GPU       & NVIDIA RTX 5070 (16~GB)    & Mali-G52 MP8 \\
RAM       & 2$\times$16~GB DDR5 5600~MHz & 8~GB LPDDR4x \\
Placa-mãe & Asus TUF Gaming B650 Plus WiFi & — \\
Fonte     & MSI A750GF (750~W, 80+ Gold) & — \\
Rede      & Gigabit Ethernet           & Gigabit Ethernet \\
\bottomrule
\end{tabular}
\end{table}

O treinamento aproveitou a RTX~5070 para aceleração CUDA do LightGBM. Os modelos
treinados foram copiados ao VIM~4 para inferência em produção.

\textbf{Ferramentas de ataque.} O \texttt{attack\_orchestrator.py} executou 8
categorias de ataque em sequência usando as seguintes ferramentas:

\begin{itemize}
  \item \textbf{nmap} — varredura de rede: SYN scan (\texttt{-sS}), UDP scan
    dos 100 top ports (\texttt{-sU}), varredura agressiva (\texttt{-A});
  \item \textbf{masscan} — varredura SYN de alta velocidade;
  \item \textbf{hping3} — geração de inundação: SYN flood (porta 80), UDP flood
    (porta 53), ICMP flood; DDoS com fonte aleatória (\texttt{--rand-source})
    nas portas 443 e 80;
  \item \textbf{medusa} — força bruta de credenciais: SSH (root), HTTP (admin),
    Telnet (root);
  \item \textbf{nikto} — varredura de vulnerabilidades web;
  \item \textbf{gobuster} — enumeração de diretórios/caminhos HTTP;
  \item \textbf{curl} — inundação HTTP (100 requisições em série);
  \item \textbf{arpspoof} — envenenamento ARP (ataque MITM);
  \item \textbf{nc} (netcat) — beacon C2 simulado via TCP porta~4444
    (100 conexões) e canal SSH reverso de longa duração.
\end{itemize}

O orquestrador executou 8 tipos de ataque em sequência contra
\texttt{192.168.100.5} (VIM~4), com sincronização confirmada via \textit{offset}
de 3~horas BRT$\rightarrow$UTC. O limiar de detecção binária foi $P_1 \geq 0{,}9$. Fluxos capturados pelo
\textit{netflower} com \texttt{idle\_timeout} = 30~s e \texttt{flow\_timeout}
= 120~s. A captura PCAP no PC atacante não foi possível por ausência de
\texttt{CAP\_NET\_RAW} no processo do tcpdump.

\subsection{IDS Binário (Fase 1)}

A granularidade da avaliação é \textbf{por segundo}: cada segundo é
classificado como ``ataque detectado'' (alerta emitido) ou não, com
\textit{ground truth} derivado das janelas de tempo do orquestrador.

\begin{table}[ht]
\centering
\caption{Matriz de confusão do IDS binário — granularidade por segundo}
\label{tab:bin_cm}
\begin{tabular}{@{}llrr@{}}
\toprule
& & \multicolumn{2}{c}{\textbf{Predição}} \\
\cmidrule{3-4}
& & \textbf{Ataque} & \textbf{Benigno} \\ \midrule
\multirow{2}{*}{\textbf{Ground truth}}
  & Ataque & 224 & 62 \\
  & Benigno &   1 & 141 \\ \bottomrule
\end{tabular}
\end{table}

\begin{table}[ht]
\centering
\caption{Métricas do IDS binário na validação ao vivo (limiar = 0,9)}
\label{tab:bin_metrics}
\begin{tabular}{@{}lr@{}}
\toprule
\textbf{Métrica} & \textbf{Valor} \\ \midrule
Acurácia          & \textbf{85,28\%} \\
Precisão          & \textbf{99,56\%} \\
Revocação (TPR)   & \textbf{78,32\%} \\
F1-score          & \textbf{87,67\%} \\
FPR               & \textbf{0,70\%} \\
FNR               & 21,68\% \\ \bottomrule
\end{tabular}
\end{table}

Os 62 falsos negativos (FN) têm origem em duas causas identificadas:
\textit{(i)} 43~s de lacuna durante o \textit{scan} UDP do nmap (\textit{recon}):
o \texttt{nmap -sU} gera poucos pacotes por porta; com \texttt{idle\_timeout}
= 30~s, muitos fluxos UDP não acumulam features suficientes para emissão
durante o \textit{scan}; \textit{(ii)} 19~s do \textit{idle slack}
pós-malware sem tráfego de ataque ativo. Os 224 FN de recon são do warmup,
não de falha do modelo — os demais tipos (DoS/DDoS, força bruta, spoofing)
atingem \textbf{100\%} de taxa de detecção por segundo.

\begin{table}[ht]
\centering
\caption{Taxa de detecção por tipo de ataque — IDS binário (granularidade por segundo)}
\label{tab:detection_rate}
\begin{tabular}{@{}lrrr@{}}
\toprule
\textbf{Ataque} & \textbf{Segundos de ataque} & \textbf{Segundos alertados} & \textbf{Taxa} \\ \midrule
recon           & 74  & 31  & 41,9\% \\
dos (incl.\ ddos) & 153 & 153 & \textbf{100,0\%} \\
bruteforce      & 16  & 16  & \textbf{100,0\%} \\
web             &  1  &  1  & \textbf{100,0\%}$^{\dagger}$ \\
mitm            &  1  &  1  & \textbf{100,0\%}$^{\dagger}$ \\
spoofing        &  3  &  3  & \textbf{100,0\%} \\
malware         & 38  & 19  & 50,0\%$^{*}$ \\ \bottomrule
\multicolumn{4}{l}{\small $^{*}$ janela inclui 30~s de \textit{idle slack};
                   os 8~s do ataque em si: 100\%.} \\
\multicolumn{4}{l}{\small $^{\dagger}$ janela de 1~s; a ferramenta concluiu em
                   $<$1~s sem serviço-alvo ativo no VIM~4 — o alerta provém de
                   fluxos DoS residuais, não de tráfego \textit{web}/ARP
                   representativo.}
\end{tabular}
\end{table}

\subsection{IDS Multiclasse (Fases 1+2)}

\begin{table}[ht]
\centering
\caption{Classes validadas — IDS multiclasse na validação ao vivo (taxonomia unificada ddos$\rightarrow$dos)}
\label{tab:multi_metrics}
\begin{tabular}{@{}lrrrr@{}}
\toprule
\textbf{Classe} & \textbf{TP} & \textbf{FP} & \textbf{FN} & \textbf{F1} \\ \midrule
recon             &   981 &     0 &   634 & \textbf{0,756} \\
dos (incl.\ ddos) & 7.033 & 1.838 &     0 & \textbf{0,884} \\ \midrule
\textbf{Acurácia geral (tipo)} & \multicolumn{4}{r}{\textbf{81,34\%}} \\
\bottomrule
\end{tabular}
\end{table}

A Fase~2 classificou corretamente \texttt{dos} (F1 = 0,884) e \texttt{recon}
(F1 = 0,756). As demais classes não puderam ser avaliadas nesta sessão por dois
motivos distintos:

\textit{(i)} \textbf{Artefato experimental} — \texttt{web} e \texttt{mitm}
dependem de serviços ativos no nó alvo (servidor HTTP para nikto/gobuster,
hosts adicionais para arpspoof). Sem esses serviços disponíveis no VIM~4, as
ferramentas completaram em menos de 1~s sem gerar tráfego representativo das
classes; a avaliação dessas categorias fica pendente para uma sessão com o
ambiente devidamente configurado.

\textit{(ii)} \textbf{Limitação de features} — \texttt{bruteforce},
\texttt{spoofing} e \texttt{malware} geraram tráfego real, mas seus padrões de
fluxo (curta duração, alta taxa de pacotes) são indistinguíveis de inundação
DoS nas 55 features por fluxo disponíveis. Superar essa limitação requer
features adicionais fora do espaço de fluxo puro, como contadores de falha de
autenticação (força bruta) ou diversidade de IPs de origem (DDoS/spoofing),
que exigem estado agregado entre fluxos.

\subsection{Desvio de Distribuição e Artefatos de Implantação}

Dois fenômenos observados em sessões anteriores de validação são relevantes
para o entendimento dos resultados e para reprodução do sistema:

\textbf{Desvio de distribuição no tráfego benigno.}
Os datasets CIC utilizam tráfego benigno scripted com padrões temporais
regulares, enquanto tráfego interativo real (SSH, HTTP de curta duração)
apresenta características distintas que o modelo classificava como ataque.
As duas features de maior ganho são as principais responsáveis: \texttt{flow\_iat\_min}
nos dados de treino benigno está na faixa de milissegundos, enquanto em bursts
TCP reais pode ser de 1--10~µs — qualquer burst no ambiente real é portanto
anômalo ao modelo. Analogamente, fluxos curtos sobre rede local rápida
produzem \texttt{flow\_byts\_s} aparente muito acima do observado no treino
benigno. A solução de médio prazo é complementar o dataset com tráfego benigno
real capturado no ambiente de implantação.

\textbf{Artefato FIN/ACK do netflower.}
Cada encerramento de conexão TCP gera um segundo micro-fluxo com
$\texttt{flow\_duration} = 0$, $\texttt{tot\_fwd\_pkts} = 1$,
$\texttt{totlen\_fwd\_pkts} = 52$~bytes e apenas flags FIN + ACK — sem
correspondente nos dados de treino benigno, o modelo atribui sistematicamente
$\sim$84\% de confiança de ataque a esses micro-fluxos. A mitigação é uma
guarda de pacote mínimo (\texttt{tot\_fwd\_pkts} $\geq 2$) no script de IDS,
descartando micro-fluxos antes da inferência.

\subsection{Comparação entre IDS Binário e IDS Multiclasse}

A Tabela~\ref{tab:comparison} consolida a comparação entre os dois pipelines
testados no VIM~4.

\begin{table}[ht]
\centering
\caption{Comparação entre dois experimentos independentes no VIM~4: pipeline apenas binário (Fase~1) e pipeline hierárquico completo (Fases~1+2)}
\label{tab:comparison}
\begin{tabular}{@{}lrr@{}}
\toprule
\textbf{Métrica} & \shortstack[r]{\textbf{IDS Binário}\\ \textbf{(só Fase~1)}}
                 & \shortstack[r]{\textbf{IDS Multiclasse}\\ \textbf{(Fases~1+2)}} \\ \midrule
\multicolumn{3}{l}{\textit{Detecção binária (Fase 1)}} \\
\quad Acurácia          & \textbf{85,28\%} & 83,89\%$^{a}$ \\
\quad Precisão          & 99,56\%          & \textbf{100,00\%$^{a}$} \\
\quad Revocação         & 78,32\%          & \textbf{83,89\%$^{a}$} \\
\quad F1-score          & 87,67\%          & \textbf{91,25\%$^{a}$} \\
\quad FPR               & 0,70\%           & N/A$^{a}$ \\ \midrule
\multicolumn{3}{l}{\textit{Classificação de tipo (Fase 2)}} \\
\quad Acurácia (tipo)   & —               & \textbf{81,34\%} \\
\quad dos F1            & —               & \textbf{0,884} \\
\quad recon F1          & —               & \textbf{0,756} \\
\midrule
\multicolumn{3}{l}{\textit{Recursos (VIM 4, durante ataques)}} \\
\quad CPU média         & 22,8\%          & 22,7\% \\
\quad CPU máxima        & 32,0\%          & 30,5\% \\
\quad RAM média         & 1.004 MB        & 1.057 MB (+53 MB) \\
\quad RAM máxima        & 1.041 MB        & 1.093 MB (+52 MB) \\ \midrule
\multicolumn{3}{l}{\textit{Throughput}} \\
\quad Alertas durante ataques & 18.482  & 9.871 \\
\quad Taxa                    & \textbf{64,8 alertas/s} & 36,0 alertas/s \\ \midrule
\multicolumn{3}{l}{\textit{Energia estimada$^{b}$ (durante ataques)}} \\
\quad Potência média          & 3,58~W & 3,78~W \\
\quad Potência média (idle)   & 2,02~W & 2,03~W \\
\bottomrule
\end{tabular}
\end{table}

\noindent$^{a}$ O log multiclasse encerra dentro da janela de ataque (sem período
benigno pós-ataque observável); FPR não computável; os valores de revocação e
F1 não incluem o \textit{idle slack} pós-malware.

\noindent$^{b}$ \textbf{Estimativa}, não medição direta — o VIM~4 não expõe
interface RAPL. Ver achado~(5) para o método de cálculo e suas limitações.

\textbf{Principais achados da comparação:}

\textbf{(1) Overhead do P2 é praticamente somente RAM.}
Durante ataques, o uso de CPU é idêntico nos dois pipelines ($\approx$22--23\%).
O custo operacional do segundo modelo é $+52$~MB de RAM ($\sim$5\% sobre a
linha de base do binário) e ausência de overhead de CPU — o modelo P2 executa
a inferência em série após P1, mas sua latência é amortizada pelo
\textit{flow tracking} que já consome o tempo dominante.

\textbf{(2) Throughput de alertas é 1,8$\times$ maior no binário.}
O IDS binário emite alerta para todo fluxo com $P_1 \geq 0{,}9$; o multiclasse
filtra adicionalmente pelo limiar P2, reduzindo o volume de alertas brutos.
Em cenários onde o operador precisa de alta frequência de alertas brutos
(e.g., para alimentar um SIEM), o binário é preferível; onde se precisa de
rótulo de tipo para triagem, o multiclasse adiciona valor sem custo de CPU.

\textbf{(3) Detecção binária equivalente.}
Como ambos usam o mesmo modelo P1, a diferença observada nas métricas binárias
decorre da janela de observação distinta, não do modelo em si. O FNR de 21,68\%
do IDS binário é causado integralmente pela lacuna do nmap UDP e pelo
\textit{idle slack} do malware — não por fluxos que o modelo deixou passar.

\textbf{(4) Beacon C2 de longa duração como caso-limite.}
O canal SSH reverso ($\sim$240~s, 980~B/s) foi detectado como ataque pelo IDS
binário (96\% de confiança). O multiclasse classificou o fluxo com baixa
confiança (melhor palpite: \textit{spoofing}, 38,6\%), evidenciando que esse
padrão — baixa taxa, longa duração, direção inversa — é inconsistente com
qualquer classe de inundação conhecida no espaço de features disponível.

\textbf{(5) Energia: estimativa, não medição direta.} O VIM~4 não expõe
interface RAPL, então a energia não pôde ser medida diretamente (campo
\texttt{power\_w} do log permaneceu vazio nas duas sessões). Em vez disso, a
potência foi \textbf{estimada} por um modelo linear de utilização de CPU:

\begin{equation}
P(t) = P_{idle} + (P_{max} - P_{idle}) \times \text{CPU\%}(t)
\label{eq:power}
\end{equation}

com $P_{idle} = 2$~W e $P_{max} = 11$~W assumidos para o SoC ARM
\textit{big.LITTLE} do VIM~4 a partir de literatura/datasheet — não
calibrados por medição em campo. A potência da Equação~\ref{eq:power} é
integrada sobre a série temporal de CPU amostrada a cada
$\sim$3~s (\texttt{SYS\_SNAPSHOT}) durante toda a sessão, e cada intervalo
entre amostras consecutivas é rotulado como ``durante ataque'' ou ``idle''
usando as mesmas janelas de \textit{ground truth} do orquestrador (a mesma
lógica de \texttt{get\_true\_label} já usada para as métricas de detecção).
O cálculo foi implementado em \texttt{scripts/ids\_metrics.py}, que reporta
potência média da sessão, durante ataques e em idle, e energia total em
joules/Wh. Como a magnitude absoluta depende da calibração assumida de
$P_{idle}$/$P_{max}$, a comparação mais confiável é a relativa: o consumo
durante ataques é praticamente o mesmo nos dois pipelines ($\approx$3,6--3,8~W),
consistente com o achado~(1) de que o custo do P2 é de RAM, não de CPU.

%─────────────────────────────────────────────────────────────────────────────
\section{Conclusão}

Este trabalho apresentou um IDS hierárquico de duas fases para ambientes IoT
em redes 6G com MEC, cobrindo desde a construção do pipeline de dados (PCAPs
$\rightarrow$ SQLite $\rightarrow$ modelos) até a implantação e validação em
produção no VIM~4. Os principais resultados são:

\begin{itemize}
  \item \textbf{Fase 1 (offline):} acurácia de 80\%, revocação de 96\% na
    classe de ataque, F2 = 0,886 ao limiar ótimo de 0,157.
  \item \textbf{IDS Binário ao vivo (limiar = 0,9):} acurácia de 85,3\%,
    precisão de 99,6\%, F1 de 87,7\%, FPR de 0,7\%. Taxa de detecção de
    100\% para DoS, força bruta, spoofing e MITM. FNR de 21,7\%, causado
    integralmente por limitações do extrator de fluxos (lacuna do scan UDP).
  \item \textbf{IDS Multiclasse ao vivo:} acurácia de classificação de tipo
    de 81,3\%, com F1 de 0,884 para DoS/DDoS e 0,756 para reconhecimento.
    Overhead de apenas 52~MB de RAM sobre o binário, sem overhead de CPU.
\end{itemize}

As principais limitações identificadas são: \textit{(i)} classes \texttt{web}
e \texttt{mitm} não puderam ser avaliadas ao vivo por ausência de serviços
ativos no VIM~4 durante o experimento; \textit{(ii)} classes \texttt{bruteforce},
\texttt{spoofing} e \texttt{malware} produzem padrões de fluxo indistinguíveis
de DoS no espaço de 55 features disponíveis, exigindo features de estado
agregado; \textit{(iii)} lacuna de detecção do \textit{recon} durante o scan
UDP do nmap, causada pelo \texttt{idle\_timeout} do \textit{netflower};
\textit{(iv)} captura PCAP no PC atacante ainda inoperante por falta de
\texttt{CAP\_NET\_RAW}.

%─────────────────────────────────────────────────────────────────────────────
\section{Trabalhos Futuros}

Como trabalhos futuros propõe-se: \textit{(i)} repetição da validação ao vivo
com servidor HTTP e hosts ARP ativos no VIM~4, permitindo avaliar as classes
\texttt{web} e \texttt{mitm}; \textit{(ii)} adição de features discriminativas
para força bruta (contadores de falha de autenticação) e spoofing/DDoS
(diversidade de IPs de origem); \textit{(iii)} redução do \texttt{idle\_timeout}
para captura de ataques de curta duração; \textit{(iv)} ativação do monitoramento
de energia no VIM~4 para quantificar o consumo real por fase.

%─────────────────────────────────────────────────────────────────────────────
\begin{thebibliography}{9}

\bibitem{akiba:19}
AKIBA, T.; SANO, S.; YANASE, T.; OHTA, T.; KOYAMA, M. {Optuna}: a
next-generation hyperparameter optimization framework. In:
\textit{Proceedings of the ACM SIGKDD International Conference on Knowledge
Discovery and Data Mining}, p. 2623--2631, Anchorage, 2019. ACM.

\bibitem{ke:17}
KE, G.; MENG, Q.; FINLEY, T.; WANG, T.; CHEN, W.; MA, W.; YE, Q.; LIU, T.-Y.
{LightGBM}: a highly efficient gradient boosting decision tree. In:
\textit{Advances in Neural Information Processing Systems}, v. 30,
p. 3146--3154. Curran Associates, 2017.

\bibitem{lazarevic:05}
LAZAREVIC, A.; ERT{\"O}Z, L.; KUMAR, V.; OZGUR, A.; SRIVASTAVA, J. Intrusion
detection: a survey. In: \textit{Managing Cyber Threats: Issues, Approaches
and Challenges}, p. 19--78. Springer, New York, 2005.

\bibitem{sharafaldin:18}
SHARAFALDIN, I.; LASHKARI, A. H.; GHORBANI, A. A. Toward generating a new
intrusion detection dataset and intrusion traffic characterization. In:
\textit{International Conference on Information Systems Security and Privacy},
p. 108--116, Funchal, 2018. SciTePress.

\bibitem{you:21}
YOU, X. et al. Towards 6G wireless communication networks: vision, enabling
technologies, and new paradigm shifts. \textit{Science China Information
Sciences}, v. 64, n. 1, p. 110301, 2021.

\bibitem{zolanvari:19}
ZOLANVARI, M.; TEIXEIRA, M. A.; GUPTA, L.; KHAN, K. M.; JAIN, R. Machine
learning-based network vulnerability analysis of industrial {Internet of
Things}. \textit{IEEE Internet of Things Journal}, v. 6, n. 4,
p. 6822--6834, 2019.

\end{thebibliography}

\end{document}
